\chapter{The differentiable chain}
\label{ch:differentiable-chain}

This chapter describes how the package makes a multireference energy
differentiable, and records the numerical facts that shaped the
design.  It is the software-side companion to the methods paper; the
derivations live there, the implementation constraints live here.

\section{Parameterization}

Orbital variation enters through a rotation parameter vector
$\theta$: the occupied--virtual mixing is generated by the
non-redundant pair list of \texttt{orblib.\allowbreak orbital\_\allowbreak pairs}, and
\texttt{rotation.\allowbreak block\_\allowbreak rotation} turns $\theta$ into the rotation
matrix applied to the reference orbitals.  Three pair families are
included: core--active, core--virtual, and active--virtual, for every
orbital outside the active set.

The active--active rotations are omitted deliberately: they are
redundant against the CI step, which is complete inside the active
space.  The core--virtual family is \emph{not} omittable.  A
parameterization that drops it stops at a point that is not
stationary; the defect was measured at the level of a few hundredths
of a kcal/mol on a test system --- small, systematic, and invisible
without an independent reference.  The pair list is the contract; do
not prune it for speed.

\section{The two factory families, and why}

Every differentiable entry point comes in two forms.

The \emph{joint} factories (\texttt{make\_nevpt2\_joint\_fn} and its
state-averaged sibling) take the molecule and $\theta$ and run the
reference SCF inside the call.  They are the geometry-differentiable
path: the nuclear coordinates enter through the integrals, and JAX
differentiates the whole composition.

The \emph{theta-only} factories (\texttt{make\_nevpt2\_theta\_fn},
\texttt{make\_nevpt2\_sa\_theta\_fn}) compute the reference SCF and the
integrals once, at construction, and return a callable of $\theta$
alone.  They exist for one hard reason: this build of pyscfad cannot
trace its SCF inside \texttt{jax.jit}, and the orbital optimizer wraps
its objective in \texttt{jit}.  A joint factory called from inside
the optimizer raises a tracing error; the theta-only factory is the
supported path there.  The two families funnel through one shared
component builder, so they cannot drift apart numerically --- the test
suite pins them against each other at fixed geometry.

\section{Inside the chain}

A single evaluation composes four numerically delicate pieces.

\textbf{The MO integral transform runs in four contraction steps.}
The four-operand form of the two-electron transform is numerically
weaker on XLA at large sizes, and the discrepancy is not cosmetic: on
an $96^4$ integral array the one-shot form showed absolute errors
around $10^{-6}$, which the near-degenerate gaps of a linear
molecule's generalized Fock matrix amplify into a completely
different eigenvector basis inside the degenerate subspace.  The
energy is insensitive to that; the theta gradient differed by eight
percent between the two contraction orders.  The stepwise form is
therefore the implementation, not an optimization.

\textbf{Canonicalization is differentiable and state-aware.}  The
NEVPT2 denominators are defined on the canonical (natural-orbital,
sub-Fock-diagonalized) set, so the contraction runs on canonicalized
orbitals produced inside the differentiated path.  The state-averaged
variant rotates with the average density of the state-averaged
solution; at one root it reproduces the ground-state canonicalization
bit for bit, which is how the two are pinned against each other.

\textbf{Per-state energies are read out of the eigensolver's
spectrum.}  The state-averaged solver returns the mean of the lowest
roots; the per-state energies are recovered from the running means by
the sliding-mean inversion $E_k = k\,M_k - (k-1)\,M_{k-1}$.  The
number of roots is silently clamped by the size of the determinant
space, so the inversion is performed over the effective root count,
not the requested one.

\textbf{Degeneracies are handled at the decomposition level.}  An
eigensolve with degenerate eigenvalues has no unique derivative basis;
the kernels use degeneracy-safe decompositions throughout, and the
choice of cluster tolerance is a protocol decision with its own
limitation entry (Chapter~\ref{ch:limitations}).

\section{What is verified, and against what}

The chain's verification record is deliberately built on invariance
and cross-implementation checks rather than on agreement with any
single reference: the CASSCF gradient is validated against invariance
conditions and against an independent derivative-free route; the
NEVPT2 energies against an independent implementation; the
state-averaged addition against the single-root path at one root.
Chapter~\ref{ch:verification} lists the records with their measured
deviations.
